\documentclass[11pt,letterpaper]{article}
\usepackage[T1]{fontenc}
\usepackage{lmodern,microtype,amsmath,amssymb,amsthm,mathtools}
\usepackage[margin=.88in,headheight=14pt]{geometry}
\usepackage{booktabs,longtable,array,enumitem,xcolor,hyperref,xurl,fancyhdr}
\definecolor{navy}{HTML}{17324D}
\hypersetup{colorlinks=true,linkcolor=navy,citecolor=navy,urlcolor=navy,
pdftitle={One modulus for all binomial digits: direct quartic consistency equations},
pdfauthor={OpenAI ChatGPT, prepared for Vladimir Reshetnikov}}
\setlength{\parindent}{0pt}\setlength{\parskip}{5.5pt}
\setlength{\emergencystretch}{2em}
\setlist{itemsep=3pt,topsep=5pt}
\pagestyle{fancy}\fancyhf{}
\fancyhead[L]{\small Direct quartic consistency equations}
\fancyhead[R]{\small 24 October 2026}\fancyfoot[C]{\thepage}
\newtheorem{theorem}{Theorem}[section]
\newtheorem{lemma}[theorem]{Lemma}
\newtheorem{proposition}[theorem]{Proposition}
\newtheorem{corollary}[theorem]{Corollary}
\theoremstyle{remark}\newtheorem{remark}[theorem]{Remark}
\newcommand{\N}{\mathbb N}\newcommand{\Z}{\mathbb Z}
\newcommand{\Con}{\operatorname{Con}}\newcommand{\ZFC}{\mathrm{ZFC}}
\newcommand{\Refute}{\operatorname{Ref}}\newcommand{\Ctwo}{\mathsf{C222}}
\newcommand{\sha}{\texttt{45b399ec51dfbecb891dd51d8a9e4b6b7c361ac6}}
\newcommand{\Psmall}{P_{332}}\newcommand{\Psparse}{P_{333}}
\title{\color{navy}One modulus for all binomial digits\\[5pt]
\large Direct quartic equations for syntactic consistency of ZFC\\
\normalsize A 332-variable equation and a sparse 333-variable companion}
\author{Prepared for Vladimir Reshetnikov and the ProveIt project\\
\small Mathematical analysis and artifact preparation with OpenAI ChatGPT}
\date{24 October 2026}
\begin{document}
\maketitle
\begin{abstract}
We construct two fixed integer-coefficient quartics whose lack of
nonnegative integer zeros is provably equivalent in ZFC to syntactic
consistency of ZFC. The first has 332 variables, 263 squared quadratic
residuals, 11,823 collected monomials, and coefficient height 85,243,454.
The second has 333 variables, 264 residuals, 4,357 monomials, and height
209,508,384. The newest inspected direct-construction baseline has 349
variables, 268 residuals, 12,768 monomials, and height 86,915,503.

The main change replaces complete binomial-numerator powers by Pell
residues in a common modulus. A new endpoint argument shows, before
interpreting a factorial, that an already required quotient-column
modulus dominates every digit modulus. Thirteen of nineteen exact power
requests disappear. A second construction arranges the free certificate
coordinates to have even indices; it removes the small-ordinate branch
without squaring the requested power outputs. This gives a much sparser
quartic with one additional variable.

The archive contains every quadratic residual, every expanded coefficient,
complete printable equations, portable generators, independent finite
checkers, and explicit ordinary proofs. The logical source is the existing
calibrated numerical calculus C222. No K215 execution, native coefficient
tuple, or Jones universal polynomial is used. The work improves a specific
small-coefficient quartic frontier. It does not establish a global simplest
equation, a minimum number of variables, or a new end-to-end proof-assistant
verification.
\end{abstract}
\newpage
\tableofcontents
\newpage

\section{The result and what ``simplest'' means here}

The comparison snapshot is ProveIt commit \sha. We inspected both the
active research and all eight archives then in \texttt{docs/incoming}.
This matters: the newest incoming direct equation already improves the
active review's 399-natural-variable point to 349 variables. The baseline
used below is that newer equation, not an earlier published count
\cite{repo,round14,active}.

Write \(\N=\{0,1,2,\ldots\}\). An equation in positive integers and an
equation in unrestricted integers are different optimization problems;
all counts in the following table use \emph{nonnegative integers}.
Coefficient height means the maximum absolute coefficient after collecting
the full sum of squares, not the largest coefficient in one residual.

\begin{center}
\begin{tabular}{@{}lrrrrr@{}}\toprule
Equation & Variables & Rows & Degree & Terms & Height\\\midrule
Newest direct parent &349&268&4&12,768&86,915,503\\
\(\Psmall\), this work &332&263&4&11,823&85,243,454\\
\(\Psparse\), this work &333&264&4&4,357&209,508,384\\\bottomrule
\end{tabular}
\end{center}

The 332-variable form saves 17 variables, five residuals, 945 monomials,
and 1,672,049 units of coefficient height. These are improvements of
approximately 4.9\%, 1.9\%, 7.4\%, and 1.9\%, respectively, with degree
unchanged. The sparse companion saves 16 variables and 8,411 monomials
relative to the parent: about 65.9\% of its expanded support disappears.
Its larger coefficient height is an explicit tradeoff. Compared with
\(\Psmall\), it uses one extra variable and 7,466 fewer monomials.

These equations do not supersede every other branch. The repository has
nine-positive-variable forms, much smaller variable counts at higher
degree, and a 37-integer quartic whose fixed coefficient tuple is very
large. This work targets a different measure: an ordinary quartic with
small literal coefficients, obtained directly from proof certificates and
with no large machine-specific numeral. There is no established total
ordering of these measures, and no global optimality theorem is claimed.

\begin{theorem}[The two delivered consistency equations]
\label{thm:main}
Let \(\Psmall\in\Z[z_0,\ldots,z_{331}]\) and
\(\Psparse\in\Z[z_0,\ldots,z_{332}]\) be the literal polynomials
specified in Section~\ref{sec:literal} and in the two accompanying complete
equation documents. Composing the ordinary arithmetic proofs in this
report with the pinned C222/refutation calibration gives
\[
 \ZFC\vdash
 \bigl[\neg\exists\boldsymbol z\in\N^{332}\;
                  \Psmall(\boldsymbol z)=0\bigr]
       \longleftrightarrow\Con(\ZFC),
\]
and the same statement for \(\Psparse\) with 333 variables.
No assumption that ZFC is consistent is used in constructing the
polynomials or proving either equivalence.
\end{theorem}

Here \(\Con(\ZFC)\) means absence of a finite syntactic derivation of
falsity from the effective ZFC axiom predicate used by ProveIt. The theorem
is an ordinary mathematical derivability claim. The delivered finite
checks verify the exact polynomials and algebraic transformations. They
do not constitute a new Lean, Rocq, or encoded object-language ZFC proof
of the entire theorem. The inherited calibration is a substantive,
explicit dependency, not a coefficient checksum.

\section{A fixed polynomial, rather than an unspecified representation}
\label{sec:literal}

The recommended equation is exactly
\begin{equation}
 \boxed{\quad\Psmall(z_0,\ldots,z_{331})
       =\sum_{j=0}^{262}Q_j(z_0,\ldots,z_{331})^2=0,
       \qquad\boldsymbol z\in\N^{332}.\quad}
 \label{eq:fixed}
\end{equation}
Every \(Q_j\) is printed in \texttt{equation-332.pdf}, generated from
\path{equations/natural332/natural/equation.json}. The full collected
polynomial, containing all 11,823 integer coefficients, is
\texttt{expanded\_polynomial.json} in that same directory. The sparse
companion is
\[
 \Psparse(z_0,\ldots,z_{332})=\sum_{j=0}^{263}\widetilde Q_j(
 \boldsymbol z)^2;
\]
its entire residual list is printed in \texttt{equation-333.pdf} and
stored under \texttt{equations/natural333/}.
In the complete equation documents, the residuals of each polynomial
are denoted simply by \(R_j\), with the same order and coefficients.

For example, the first three residuals of \(\Psmall\), in the delivered
coordinate order, are
\begin{align*}
 Q_0&=z_3-z_0^2-24z_0-144,\\
 Q_1&=z_4-z_2z_3,\\
 Q_2&=z_5+z_7-z_2-z_2z_6-z_2z_7.
\end{align*}
Thus the horizon begins as \(t=z_0+12\), its square is stored in
\(z_3\), and the early rows include the factorial ceiling identity.
The equation is not just these illustrative rows: all 263 residuals
are part of definition~\eqref{eq:fixed}.

A sparse term \texttt{[c,[i,j,...]]} in the JSON means
\(c z_i z_j\cdots\); repeated indices mean repeated factors and an
empty index list denotes a constant. All indices are zero-based, and
the coordinate order is fixed. There are no unspecified constants,
variable exponents, factorials, binomial coefficients, floor functions,
or unbounded products left in either final polynomial. Those functions
appear below only to explain the meanings of finite polynomial gadgets.

Every residual has degree at most two. Squaring therefore gives degree
at most four. The independent full expansion finds nonzero quartic
coefficients, proving that the degree is exactly four on the actual
literal coefficients. No random evaluation is used to infer the degree.

\subsection{Other domains}

The same 332-variable equation over strictly positive integers is
\(\Psmall(y_0-1,\ldots,y_{331}-1)=0\). This is a bijective change of
domain, although its expanded coefficient height has not been optimized.
For an explicit unrestricted-integer equation, put
\[
 \theta(a,b,c)=a^2+b^2+c(c+1).
\]
This polynomial is always nonnegative on integers and represents every
nonnegative integer. Indeed the three-square theorem applied to
\(4n+1\) gives one odd and two even squares; rearranging gives
\(n=\theta(a,b,c)\). Consequently
\[
 \Psmall\bigl(\theta(a_0,b_0,c_0),\ldots,
                    \theta(a_{331},b_{331},c_{331})\bigr)=0
\]
is a completely specified unrestricted-integer companion in 996 variables,
of degree eight. It is a domain conversion, not a claim to improve the
repository's dedicated integer constructions. The new quartic counts
always refer to \(\N\).

\section{Logical source: a finite certificate of a refutation}

The construction uses the calibrated numerical calculus \(\Ctwo\)
already present in ProveIt. Its formula constructors begin with
\[
 I(a,b)=2(a+b)(a+b+1)+4b.
\]
The relation, universal-quantifier, and equality constructors add 1, 2,
and 3, respectively. The code 0 is truth and the code 1 is falsity.
There are nine fixed axiom/schema constructor polynomials, together with
truth, modus ponens, and generalization. The bundle schema has the side
conditions \(c\ne a,b\), represented by
\[
 ((c-a)(c-b))^2=(d+1)^2,\qquad d\in\N.
\]
This equation is equivalent to both inequalities: when the product is
nonzero, take \(d=|(c-a)(c-b)|-1\).

The exact constructor definitions and rule declarations are retained in
the vendored \texttt{Code254.lean} and \texttt{KCode222.lean} files and
in the literal source JSON. The relevant calibration interfaces are
\begin{align*}
 \Ctwo(1)&\longleftrightarrow
       \Refute_{\mathrm{ZF}}(\bot),\\
 \Refute_{\mathrm{ZF}}(\bot)&\longleftrightarrow
       \Refute_{\mathrm{ZFC}}(\bot).
\end{align*}
Their exact names are \path{KEndpoint.c222_iff_zf_refutation}
and \path{ConstructibleChoice.zf_refutation_iff_zfc_refutation}.
The second uses the constructible-choice equiconsistency transport.
The supplied source pin and hashes identify these existing statements
\cite{calibration}. They do not involve a machine execution hypothesis.

This use of a direct proof calculus is consistent with the broader
program of arithmetic representations of provability studied by Carl
and Moroz and, more recently, Formisano, Gavazzi, and Omodeo
\cite{carlmoroz,fgo}. The specific source and optimizations here are the
ones identified in the ProveIt provenance; those papers are context,
not substitutes for this source's calibration.

\subsection{One factorial serves two purposes}

Choose a horizon \(t\ge12\), eventually the number of padded proof
steps, and let \(F=t!\). Every formula label in a local step is required
to be at most \(t\). Define its weight by
\[
 w_j=1+jF,\qquad 0\le j\le t.
\]
For distinct positive \(j,k\le t\), these weights are coprime. A common
prime divisor would divide \((j-k)F\), would not divide \(F\), and
would therefore divide \(j-k\); but that prime is at most \(t\), hence
divides \(t!\), a contradiction. Truth has the harmless weight \(w_0=1\).

A local rule with premise labels \(p,q\) and conclusion \(v\) is
encoded by
\begin{equation}
 S' w_v=S w_p w_q. \label{eq:flow}
\end{equation}
Unused premises are truth. The accumulator starts at \(S_0=w_1\)
and ends at \(S_t=1\). Since every weight is positive, its local
equalities propagate positivity of all accumulator entries. Multiplication
of~\eqref{eq:flow} and cancellation give
\[
 \prod_{i<t}w_{v_i}=w_1\prod_{i<t}w_{p_i}w_{q_i}.
\]
Prime valuations isolate each positive label and give
\begin{equation}
 \#\{i:v_i=j\}-\#\{i:p_i=j\}-\#\{i:q_i=j\}
       =[j=1]. \label{eq:balances}
\end{equation}

These balances force a derivation even if the certificate's rule multiset
contains cycles. Let \(D\) be the closure of truth under its valid rule
instances. If falsity were outside \(D\), every rule with conclusion
outside \(D\) would have a premise outside \(D\). Summing
\eqref{eq:balances} over the finitely many positive labels outside \(D\)
would give a nonpositive left side and the right side 1. Hence falsity
is in the closure. Conversely, expanding a finite proof tree backwards
and storing the product of the weights of its pending obligations gives
the required positive accumulator. Truth steps pad it to any sufficiently
large horizon.

The same factorial also gives pairwise-coprime reciprocal moduli
\[
 d_i=1+\frac{F^2}{i}\quad(1\le i\le t),\qquad
 B=\prod_{i=1}^t d_i=\binom{F^2+t}{t}.
\]
Each \(d_i\equiv1\pmod F\), so its prime divisors exceed \(t\).
A common prime divisor of \(d_i,d_j\) would divide \(i-j\), again
impossible. The divisions and product in this paragraph describe the
proof; the actual graph uses the displayed binomial identity.

\subsection{Amplified CRT and a derived label range}

Set
\[
 R=2+t+x+y+F,\qquad K=2R,\qquad N=B^K.
\]
A beta code \((x,y)\) stores the overlapping accumulator entries with
moduli \(1+(i+1)y\). Eight CRT columns store the four local rule
parameters, two accumulator coordinates, and two quotient coordinates. A ninth
representative stores the row index. Their range conditions have the form
\(N\mid\binom{a}{k}\).

The key range implication is
\begin{equation}
 p^e\mid\binom{a}{k}\quad\Longrightarrow\quad
                 a\bmod p^e<k. \label{eq:range}
\end{equation}
For \(a\ge k\), if \([a-k+1,a]\) contains no multiple of \(p^e\),
then the valuation formula for the binomial has zero terms from level
\(e\) onwards and at most \(e-1\) remaining units. That contradicts
the assumed divisibility. The case \(a<k\) is immediate. Only this
direction is used; its converse is generally false.

To recover a row, choose a prime of its reciprocal modulus, use
\eqref{eq:range} to decode all bounded coordinates, and reduce the local
nonnegative source polynomial modulo the corresponding amplified prime
power. Its absolute size is below that prime power, so its zero residue
is an actual zero. Conversely, CRT at precision \(d_i^{2K}\) gives
representatives whose falling factorials have enough valuation even after
division by \(k!\). The factor 2 in the completeness precision is
essential.

Instead of another local label-bound coordinate, define \(J\) as the
positive polynomial sum of all nine schema codes and all labels that
could occur in the chosen local rule. Require
\(N\mid\binom{J}{t+1}\). A separate majorant bounds this derived
expression. It must not be included in the upper bound on independently
chosen CRT representatives. This is the established derived-range
architecture of the newest parent \cite{round14}; both delivered equations
retain that distinction.

\section{The new dominant denominator}
\label{sec:dominant}

The old arithmetic graph evaluates complete powers \((q+1)^n\) to
extract a small binomial digit. We retain only what the digit actually
uses: the residue of that power in a suitable modulus. Package 108
already introduced the cofinal Pell method for such residues
\cite{modular108}. The new issue is to combine it with the latest
factorial architecture without paying for a new large denominator power.

The semantic endpoint equations are
\begin{align}
 x&=F+1+(y+1)u_0,\nonumber\\
 y&=F+r_0+1,\label{eq:endpoints}\\
 x&=1+[1+(t+1)y]u_t.\nonumber
\end{align}
Before its factorial meaning is known, the proposed output \(F\) occurs
in the ceiling certificate
\begin{align}
 \rho&=Ft^2,& v&=(\rho+1)^t,\nonumber\\
 v+c&=cF+r+1,&c&=r+s+1. \label{eq:ceiling}
\end{align}
Initially \(c\) is only a proposed value of
\(\binom{\rho+t}{t}\). All witnesses here are nonnegative.

\begin{lemma}[Endpoint dominance before factorial interpretation]
\label{lem:dominant}
Once the retained exact power in~\eqref{eq:ceiling} has been recovered,
the raw polynomial rows imply
\[
 F\ge1,\qquad y\ge2,\qquad
 x\ge2+(t+1)y>\max(t+1,y+1).
\]
No assertion that \(F=t!\) or that \(c\) is a binomial is needed.
\end{lemma}
\begin{proof}
The power gives \(v>0\). If \(F=0\), the last two equations of
\eqref{eq:ceiling} give \(v=-s\le0\). Thus \(F\ge1\), and
\eqref{eq:endpoints} gives \(y\ge2\) and \(x\ge2\). Its final
equation forces \(u_t\ge1\), giving the claimed lower bound for \(x\).
The strict comparisons use only \(t\ge1\) and \(y\ge2\).
\end{proof}

The radix is
\begin{equation}
 q=N(N+v+1)^K. \label{eq:radix}
\end{equation}
Every lower digit index of the smaller construction belongs to
\(\{t,t+1,x+1,y+1\}\). Exact binomials use lower index \(t\).
Retain exact powers \(q^t,q^{x+1},q^{y+1}\) and define
\(q^{t+1}=q\,q^t\). The single common modulus is
\begin{equation}
 \boxed{M=Nq^{x+1}.} \label{eq:dominantmod}
\end{equation}
This is already the modulus required for the two quotient-column range
tests. By Lemma~\ref{lem:dominant}, every range modulus \(Nq^k\)
divides \(M\); the exact-digit modulus \(q^{t+1}\) does too.
The proof uses no unvalidated factorial to establish the comparison.

The six remaining exact powers are
\[
 (Ft^2+1)^t,\quad B^K,\quad (N+v+1)^K,\quad
 q^t,\quad q^{x+1},\quad q^{y+1}.
\]
All twelve binomial numerator powers disappear. Replacing
\(q^{t+1}\) by a multiplication removes one further power request,
giving nineteen down to six. The numerator indices still have certified
Pell pairs; only their full power outputs disappear.

\section{Pell residues with one cofinal choice of parameter}

For \(A\ge2\), write
\[
 X_n(A)+Y_n(A)s=(A+s)^n,
       \qquad s^2=A^2-1.
\]
The integer sequences start at \((1,0)\), then \((A,1)\), and follow
the usual multiplication recurrence. The classification and index
properties used here are the standard Pell ingredients of the
Diophantine representation of powers \cite{pell}; the exact shared
index certificate is retained in the vendored source.

\subsection{A lower bound without fixing the parameter}

Let \(W\) be 2 plus the sum of distinct semantic atoms appearing as
bases, exponents, or proposed outputs of the six exact-power requests.
Impose
\begin{equation}
 Z=z+1,\quad
 L^2-\bigl((W+1)^2-1\bigr)(WZ)^2=1,\quad
 A=L+\sigma,\qquad z,\sigma\in\N. \label{eq:cofinal}
\end{equation}
Pell classification gives \(L=X_e(W+1)\) and \(WZ=Y_e(W+1)\).
The recurrence modulo \(W\) gives \(Y_e(W+1)\equiv e\pmod W\).
Thus \(e\ge W\), and
\begin{equation}
 A\ge L\ge(W+1)^W>W^W. \label{eq:framegrowth}
\end{equation}
Crucially, after selecting this lower bound, \emph{any} larger \(A\)
is available through \(\sigma\).

For a proposed exact power \(m=n^k\), the graph imposes
\[
 T_n=2An-n^2-1,\qquad
 X_k+nY_k=AY_k+m+T_n h.
\]
Its nonnegative modulus row excludes \(n=0\). The frame bound makes
\(T_n>m,n^k\), while the Pell recurrence gives
\(X_k+(n-A)Y_k\equiv n^k\pmod{T_n}\). The residue is therefore exact:
\(m=n^k\). This recovers the retained powers before interpreting
binomials or factorials.

\subsection{One root congruence for every numerator}

Put \(b=q+1\). Impose
\begin{equation}
 2Ab=b^2+1+Mh_0,\qquad h_0\in\N. \label{eq:root}
\end{equation}
Because \(N\mid q\), we have \(\gcd(b,M)=1\). The linear
congruence for \(A\) is soluble: \(\gcd(2b,M)\) is either 1 or 2,
and in the second case \(q\) is even, so \(b^2+1\) is even. Its
solutions form an arithmetic progression. We may choose \(A\) in that
progression above both \(L\) and \(b\), which also makes \(h_0\)
nonnegative. This is compatible with all finitely many shared index
witnesses. All genuine semantic data are chosen before this \(A\).

Equation~\eqref{eq:root} makes the substitution \(s\mapsto b-A\)
a homomorphism to \(\Z/M\Z\), because
\[
 (b-A)^2-(A^2-1)=b^2-2Ab+1\equiv0\pmod M.
\]
It follows, for every \(n\ge0\), that
\begin{equation}
 \boxed{X_n+(b-A)Y_n\equiv b^n\pmod M.} \label{eq:pellresidue}
\end{equation}
This identity is the reason an enormous numerator need not be produced.

For \(n\ge1\), the norm also gives
\[
 X_n^2-(A-1)^2Y_n^2=1+2(A-1)Y_n^2>0.
\]
Hence
\begin{equation}
 X_n+(b-A)Y_n>(b-1)Y_n>0. \label{eq:positive}
\end{equation}
The positivity is needed for completeness over \(\N\); merely proving
a congruence would not justify natural quotient witnesses.

\subsection{Extracting the actual digits}

Let \(d=q^k\). An exact binomial digit \(c\) is represented by
\begin{align}
 X_n+bY_n&=AY_n+r+dc+qd h,\nonumber\\
 d&=r+s+1,&q&=c+o+1. \label{eq:exactdigit}
\end{align}
Thus \(0\le r<d\) and \(0\le c<q\). Since \(qd\mid M\),
\eqref{eq:pellresidue} identifies \(r+dc\) as the least residue of
\((q+1)^n\) modulo \(q^{k+1}\). When the low coefficients
\(\binom nj\), \(0\le j\le k\), are smaller than \(q\), the
binomial expansion has no carries there. Unique quotient and remainder
then give \(c=\binom nk\).

For a range test only divisibility by \(N\) is needed. Replace
\eqref{eq:exactdigit} by
\begin{align}
 X_n+bY_n&=AY_n+r+Nd h,&d&=r+s+1. \label{eq:rangedigit}
\end{align}
Because \(N\mid q\), the residue of \((q+1)^n\) modulo \(Nd\)
has the form
\[
 r_0+d\left(\binom nk\bmod N\right),\qquad 0\le r_0<d.
\]
It lies below \(d\) exactly when \(N\mid\binom nk\). This proves
the desired range semantics. Conversely choose the true digits and
remainders. The positive integer in~\eqref{eq:positive} has the same
least residue, so its quotient is nonnegative. Thus every displayed
new row has witnesses in the required domain.

\section{Bootstrap order and the finite horizon constant}

The dependencies cannot be rearranged arbitrarily. In particular a
factorial estimate must not be used to validate the digit that first
establishes that factorial. The soundness order is:

\begin{enumerate}
\item Restore the raw gate system from its checked elimination certificates.
\item Recover the exact powers through the cofinal frame and shared Pell
      indices. Zero power bases are excluded by their raw modulus rows.
\item Apply Lemma~\ref{lem:dominant}; establish that every requested digit
      modulus divides \(M\), without factorial semantics.
\item Apply the common root congruence and the Pell residue identity.
\item Validate the first exact binomial digit, recover \(F=t!\), then
      validate the reciprocal product binomial and all range digits.
\item Decode bounded rows by amplified CRT, extract a C222 derivation by
      the finite Horn balance argument, and use the pinned calibration.
\end{enumerate}

For the first digit, exact powers give \(q>v^2\) and \(v=(Ft^2+1)^t\).
After \(F\ge1\) is known from raw rows, the low coefficients of
\((q+1)^{Ft^2+t}\) are bounded by \(v<q\). The ceiling identity
\eqref{eq:ceiling} then yields \(F=t!\). Next
\(F^2+t\le(Ft^2+1)^2\), so the low coefficients of the reciprocal
product digit are at most \(v^2<q\). This establishes \(B\ge2\)
and \(N\ge2^K\) before the remaining range estimates.

The parent majorants use
\[
 E=8,388,608,\qquad \ell=4,194,304.
\]
They bound the local source by \(R^E\), and the derived label
polynomial on bounded inputs by \((n+1)^\ell\). These are proof
bounds, not additional coefficients or unknowns in the final equation.
The older horizon floor 16 can be lowered to 12 because
\begin{align*}
 12!&=479,001,600>2^{28},\\
 28E&=234,881,024<2^{28},\qquad 2E<2^{28},\\
 2\cdot12!&=958,003,200>3\ell\cdot13=163,577,856.
\end{align*}
At \(r=2^{28}\) these inequalities give \(r^E<2^r\).
For larger integers \(r\),
\[
 (1+1/r)^E\le\sum_{j=0}^{E}(E/r)^j<2,
\]
so induction preserves the bound. Hence, after factorial interpretation,
\(R^E<2^R<2^K\). The last displayed factorial inequality similarly
propagates to \(2t!\ge3\ell(t+1)\), supplying the derived-digit bound.
Completeness can always pad a finite proof beyond this minimum; no
finite proof length is being bounded in advance.

\section{A sparse companion by arranging even indices}
\label{sec:parity}

The second contribution changes the certificate representation itself.
An earlier incoming method doubled every power exponent and replaced
\(m=n^k\) by \(m^2=n^{2k}\) to remove exceptional Pell-index cases
\cite{even}. Here we arrange the free source coordinates so that the
\emph{actual} exponents are already even. The power outputs remain
unsquared. This is an existence-preserving change of certificates, not
a pointwise substitution into every old witness.

The free outer coordinates become
\begin{equation}
 t=2t_0+12,\quad x=2x_0+2,\quad F=2f_0,\quad
 y=F+2r_0+2. \label{eq:evenfront}
\end{equation}
Every independently chosen CRT representative is \(2(a_0+1)\).
Its aggregate bound doubles from \(9N^2\) to \(18N^2\).
The four local parameter bounds are now strict below \(t\), the two
accumulator bounds strict below \(y\), and the two quotient bounds
strict below \(x\). Accordingly the lower range indices are
\(t,y,x\), without the old additions of 1. The derived label range is
\begin{equation}
 N\mid\binom{2J}{2t}. \label{eq:evenlabels}
\end{equation}
Its decoding gives \(2J<2t\), hence \(J<t\). All required upper and
lower Pell indices are syntactically even before any factorial or
binomial meaning is invoked. In particular \(F=2f_0\) is enforced
as a polynomial definition; parity is not inferred prematurely from
the intended factorial.

\subsection{Why the parity restrictions preserve completeness}

Choose an even padded horizon strictly above all finitely many source
parameters and label sums, then set \(F=t!\). It is even. Choose an
even \(y\), divisible by \(t!\), strictly larger than all accumulator
values and at least \(F+2\). The beta moduli \(1+(i+1)y\) are odd.
Their joint CRT solution can therefore be made positive and even by
adding a multiple of their odd product; make it large enough to obtain
the prescribed \(x=2x_0+2\). Its quotient columns are strictly below
\(x\) because the relevant positive divisors exceed 1.

The reciprocal moduli \(d_i\) are odd as well. A residue class modulo
their product at precision \(2K\), namely \(N^2\), has a positive
even representative at most \(2N^2\). Taking all nine representatives
this way establishes the aggregate bound. Multiplication by 2 in
\eqref{eq:evenlabels} is compatible with the same high-precision CRT
lifting. It is not an appeal to the false converse of
\eqref{eq:range}.

The ordinary ranges satisfy the old carry estimates with 18 in place
of 9. For the derived expression, the separate majorant gives
\(2J\le N^{3\ell+1}\). It suffices that
\[
 2t(3\ell+1)\le2F<K.
\]
At \(t=12\), its left half is
\(12(3\ell+1)=150,994,956<12!\); multiplication by successive
integers preserves the required factorial inequality. This verifies the
stronger derived-label digit budget independently of the aggregate bound.

\subsection{Even-index exclusion, with no hidden sign assumption}

For each seed index \(k\), the shared certificate has main and auxiliary
Pell equations and congruences
\begin{align*}
 X^2-(A^2-1)Y^2&=1,&S^2-(B_*^2-1)T^2&=1,\\
 S&\equiv X\pmod U,&T&\equiv k\pmod{4Y}.
\end{align*}
The common parameters satisfy
\[
 U^2-(A^2-1)V_*^2=1,\quad V_*\ge1,\quad
 B_*\equiv A\pmod U,\quad B_*\equiv1\pmod{4V_*},\quad
 Y^2\mid V_*.
\]
The literal product defining \(V_*\), together with its strict
positive definition, forces every seed \(Y\ne0\).

If \(Y=1\), the main norm gives \(X=A\). The standard abscissa
residue lemma says
\[
 X_j(A)\equiv\pm A\pmod{X_h(A)}\quad\Longrightarrow\quad j
 \text{ is odd}.
\]
One proof reduces the index modulo \(2h\) using
\((A+s)^{2h}\equiv-1\pmod{X_h(A)}\), reflects into \([0,h]\),
and uses strict increase of \(X_j\) to force the reduced index to 1.
Apply this to the auxiliary Pell pair. Its index is odd, and its odd
base \(B_*\) makes its ordinate odd. But
\(T\equiv k\pmod{4Y}\) and the even \(k\) make that ordinate even.
Contradiction. Thus \(Y\ge2\), and Pell growth gives
\(Y\ge2A>k\).

For the last inequality, the sparse construction includes every
modular-only upper index in \(W\) as well as the exact-power atoms.
Omitting those extra indices would leave an unjustified bound for the
large derived expression \(2J\). With the complete frame, the ordinary
index criterion now recovers the requested index without a separate
\(Y-k\ge0\) coordinate. All variables are natural throughout this
argument. It is not a signed-domain proof and cannot be transferred to
unrestricted integers by simply dropping the same guards.

The dominant-modulus argument strengthens under~\eqref{eq:evenfront}:
raw rows give \(F\ge2\), \(y\ge4\), and
\(x\ge2+(t+1)y>2t,y\). Hence the sparse common modulus is
\[
 M=Nq^x.
\]
Its extra denominator is obtained by \(q^{2t}=(q^t)^2\), so it too
has only six exact-power requests. The common root, positive residue,
and digit proofs then apply unchanged. The complete quantified proof
and literal parity checker are included separately in the archive.

\section{Exact verification and the proof boundary}

The finite verification has two distinct purposes. First, it binds the
mathematical formulas to the actual emitted rows. Second, it checks the
full coordinate eliminations and every final coefficient. Neither purpose
is a bounded search for a refutation or a test of consistency.

For \(\Psmall\), the independent checker restores all 914 original
quadratic rows through 496 discarded gate coordinates, checks 155 further
monic eliminations, and expands all 11,823 coefficients. The two
restorations with negative coefficients have explicit rational
sum-of-squares nonnegativity certificates; coefficient signs are not
silently treated as a positivity proof. The unchanged 441-row source
prefix is compared with the pinned parent, apart from the recorded
horizon constant. A second checker verifies all 47 new macro identities
over exact sparse polynomial rings. Deliberate changes to a restoration
coefficient and a polynomial coefficient are rejected.

For \(\Psparse\), the same independent algebraic checker handles
922 original rows, 516 gate restorations, 142 further monic contractions,
and all 4,357 final coefficients. Its separate structure check verifies
the syntactic parity, actual altered source bounds, complete frame, shared
positive ordinate, dominant modulus, and absence of deleted index-gap
rows. The ordinary parity and CRT proof remains necessary beyond these
finite checks.

The archive retains the ordinary source proofs, calibration declarations,
source identities and licensing notices. The complete literal equations
and their checkers can be used without network access. Generator replay
uses only the supplied Python modules and the standard library; LaTeX
is needed only to rebuild the PDFs. The verification receipts state their
actual scope. No test output is relabeled as a new proof-assistant
certificate.

\subsection{Why nonnegative elimination is checked explicitly}

An affine gate such as \(z=x+y\) can be deleted over \(\N\) because
its restoring value is always nonnegative. A general monic row
\(z=f(\boldsymbol x)\) cannot be deleted merely because it has
coefficient 1 at \(z\): the polynomial \(f\) might take negative
values. The supplied certificates give a coefficientwise nonnegative
formula or a verified nonnegative square decomposition for every
restoration used. Each deleted coordinate is uniquely restored, and all
retained rows are checked after substitution. This is the finite step
that carries the high-level arithmetic construction to the stated exact
variable counts.

\section{What this round establishes, and what remains open}

The completed contribution is a pair of explicit consistency equations
and their proof-preserving compiler changes. The new endpoint dominance
lemma removes a separate common-denominator power; the combined modular
construction removes thirteen exact-power requests; parity padding gives
a sparse alternative with unsquared power outputs. The repository-relative
comparison includes the latest incoming archives. The generic Pell residue
identity, CRT, coprime label weights, and even-index obstruction are
credited to their existing sources rather than claimed as newly discovered
general theorems.

Several tempting shortcuts were not used. The common Pell parameter is
chosen afresh during completeness, so the proof does not falsely retain
an old frame that bounded much larger numerator outputs. The derived
label expression is not squeezed into the bound on independent CRT
representatives. Parity is established by literal constructors before
factorial interpretation. Signed modular-only Pell pairs are not assumed
to have the orientation supplied by a removed exact-power call.

The next meaningful reductions would require a further change in the
certificate or index architecture, or a sound reduction of its local
source polynomial. Removing another positivity condition without a
replacement argument is not such a reduction. Likewise, combining two
radicals into one square can introduce nontrivial squareclass solutions;
finite searches do not justify discarding those branches. Those approaches
remain research questions, and no unproved candidate is included among
the delivered consistency equations.

This is measurable progress on a small-coefficient quartic representation
of ZFC refutability. It supplies two new explicit points for comparison;
it does not supply a lower bound establishing global optimality for
either one.

\appendix
\section{Reading and rebuilding the package}

\begin{center}
\begin{tabular}{@{}p{.36\textwidth}p{.57\textwidth}@{}}\toprule
Location & Purpose\\\midrule
\texttt{report.tex}, \texttt{report.pdf} & This research article.\\
\texttt{equation-332.tex/.pdf} & Every residual of the 332-variable polynomial, in order.\\
\texttt{equation-333.tex/.pdf} & Every residual of the sparse companion.\\
\texttt{equations/natural332/} & Raw system, exact projection, all residuals and collected coefficients.\\
\texttt{equations/natural333/} & Corresponding complete sparse-system artifacts.\\
\texttt{code/} & Portable producers and independent checkers.\\
\texttt{proofs/} & Complete ordinary transformation proofs and independent review.\\
\texttt{vendor/round14/} & Pinned prior architecture, source calculus and arithmetic dependencies.\\
\texttt{verification/} & Actual finite verification and reproduction receipts.\\
\texttt{SOURCE\_LEDGER.json} & Compared repository snapshot, archive identities and inherited source map.\\
\texttt{MANIFEST.json} & SHA-256 and byte length of every packaged file except the manifest itself.\\\bottomrule
\end{tabular}
\end{center}

The top-level \texttt{verify\_all.py} verifies the manifest and both
equations, including their independent algebraic and structural checks.
Use \texttt{python3 verify\_all.py --rebuild} to generate both equations
again in a temporary directory and compare every selected equation and
coefficient file. The PDF rebuild command is
\texttt{latexmk -pdf -interaction=nonstopmode -halt-on-error report.tex};
the two complete equation documents compile the same way. See the package
README for exact commands and the interpretation of each verification.

\begin{thebibliography}{10}\raggedright
\bibitem{repo} V. Reshetnikov and ProveIt contributors,
\emph{ProveIt}, snapshot \sha, inspected 24 October 2026.
\url{https://github.com/VladimirReshetnikov/ProveIt/tree/45b399ec51dfbecb891dd51d8a9e4b6b7c361ac6/Computability/DiophantineConsistency}.
\bibitem{round14} ProveIt incoming research package,
\emph{Powered Saturation and Factorial Carriers, Round 14},
\texttt{ConZFC\_Powered\_Saturation\_Factorial\_Carriers\_Round14\_2026-10-24.zip},
in \texttt{docs/incoming} at \cite{repo}. The used components are
\texttt{architecture/DERIVED\_LABEL\_RANGE\_PROOF.md},
\texttt{FACTORIAL\_CARRIER\_PROOF.md}, and \texttt{HYBRID\_PROOF.md}.
\bibitem{active} ProveIt, \emph{Later arrivals and current size reductions},
\texttt{PROPOSAL\_REVIEW\_LATEST.md}, at \cite{repo}.
\bibitem{modular108} ProveIt package 108,
\emph{Modular binomial certificates through a cofinal Pell frame},
\texttt{108-pell-residues/proofs/modular\_binomial.txt}, at \cite{repo}.
\bibitem{even} ProveIt incoming research package,
\emph{Weaker congruences and even Pell indices},
\texttt{ConZFC\_Carriers\_and\_Even\_Indices\_2026-10-24.zip},
especially \texttt{direct566/proofs/even\_power\_compiler.md}, at \cite{repo}.
\bibitem{calibration} ProveIt, \emph{Ordinary correctness review: direct C222 quartic, round 12},
\texttt{PROPOSAL\_REVIEW\_DIRECT\_C222\_ROUND12.md}, at \cite{repo};
and the vendored calibration declarations with their independent source
pin \texttt{08200eb10fde2db408c69dffd840d6920a569b1f}.
\bibitem{carlmoroz} M. Carl and B. Z. Moroz,
\emph{On a Diophantine representation of the predicate of provability},
Max Planck Institute preprint 2006/129.
\url{https://webdoc.sub.gwdg.de/ebook/serien/e/mpi_mathematik/2006/129.pdf}.
\bibitem{fgo} A. Formisano, I. Gavazzi, and E. G. Omodeo,
\emph{Diophantine Modeling of Provability in Algebraic Logic},
CILC 2025, CEUR Workshop Proceedings 4003, paper 13.
\url{https://ceur-ws.org/Vol-4003/paper13.pdf}.
\bibitem{pell} Lean community, \emph{Mathlib.NumberTheory.PellMatiyasevic},
Pell classification, index criteria, and representation of powers,
documentation inspected 24 October 2026.
\url{https://leanprover-community.github.io/mathlib4_docs/Mathlib/NumberTheory/PellMatiyasevic.html}.
\end{thebibliography}
\end{document}
